%ID VD_3_3_2_01
\textbf{Bước 1. Chọn sine theo biên cố định.} Trên $(0,L)$, $X_n=\sin(n\pi x/L)$, $\omega_n=an\pi/L$. Hệ số của li độ là $A_n$, của vận tốc là $B_n$; nghiệm mode là $(A_n\cos\omega_nt+B_n\sin\omega_nt/\omega_n)X_n$.

\textbf{Bước 2. Câu (i).} $a=L=1$, $f=\sin\pi x$, $g=\sin^3\pi x=(3\sin\pi x-\sin3\pi x)/4$. Trực giao cho $A_1=1$, $B_1=3/4$, $B_3=-1/4$, các hệ số khác bằng $0$. Do đó
\[u=\sin\pi x\cos\pi t+\frac3{4\pi}\sin\pi x\sin\pi t-\frac1{12\pi}\sin3\pi x\sin3\pi t.\]
Đạo hàm thời gian tại $0$ khôi phục $g$; từng mode thỏa PDE và biên. Nghiệm $C^\infty$, chu kỳ cơ bản $2$.

\textbf{Bước 3. Câu (ii).} $a=2$, $L=\pi$, dữ liệu trong sách là $f(x)=\sin(\pi x)$, không phải $\sin x$. Dùng $2\sin\pi x\sin nx=\cos((\pi-n)x)-\cos((\pi+n)x)$:
\[A_n=\frac2\pi\int_0^\pi\sin(\pi x)\sin nx\,dx
=\frac{2(-1)^nn\sin(\pi^2)}{\pi(\pi^2-n^2)},\qquad B_n=0.\]
Không có mẫu số bằng không vì $n$ nguyên. Nghiệm là $u=\sum_{n\ge1}A_n\cos2nt\sin nx$, chu kỳ $\pi$.

\textbf{Bước 4. Nêu mức trơn đúng.} $f(\pi)=\sin(\pi^2)\ne0$, trái biên $u(\pi,t)=0$. Do đó câu (ii) không có nghiệm liên tục đến góc $(\pi,0)$ với mọi dữ liệu nguyên văn. Chuỗi xác định nghiệm theo nghĩa $L^2$/ từng miền; bước nhảy của mở rộng lẻ phát sinh mặt sóng từ đầu phải. Câu (i) không có vấn đề tương thích ấy.
%ID VD_3_3_2_02
\textbf{Bước 1. Giữ mode hằng của Neumann.} Hệ hàm riêng là $X_0=1$, $X_n=\cos(n\pi x/L)$, $n\ge1$. Mode $0$ bằng $A_0+B_0t$ với hệ số là trung bình của dữ liệu.

\textbf{Bước 2. Câu (i).} $a=L=1$, $f=0$, $g=\cos^2\pi x=(1+\cos2\pi x)/2$. Vì vậy $B_0=1/2$, $B_2=1/2$, còn lại bằng $0$:
\[u=t/2+\cos2\pi x\sin2\pi t/(4\pi).\]
Đạo hàm không gian là bội của $\sin2\pi x$, bằng không ở hai đầu. Đạo hàm thời gian tại $0$ cho đúng $g$. Hạng tuyến tính khiến nghiệm không tuần hoàn dù phần dao động có chu kỳ $1$.

\textbf{Bước 3. Câu (ii).} $a=2$, $L=\pi$, $f=\sin x$, $g=0$. $A_0=2/\pi$. Tích–tổng cho
\[A_n=\frac2\pi\int_0^\pi\sin x\cos nx\,dx
=\frac{2(1+(-1)^n)}{\pi(1-n^2)}\quad(n\ne1),\qquad A_1=0.\]
Tại $n=1$ phải tích phân riêng, không dùng mẫu bằng $0$. Chỉ còn các mode chẵn:
\[u=\frac2\pi+\sum_{m\ge1}\frac4{\pi(1-4m^2)}\cos4mt\cos2mx.\]
\textbf{Bước 4. Chu kỳ và trơn.} Mode $m=1$ khác không cho chu kỳ cơ bản $\pi/2$. Dữ liệu $f'(0)=1$, $f'(\pi)=-1$ không thỏa biên Neumann; mở rộng chẵn có góc, nên không có nghiệm $C^1$ trên toàn miền qua các mặt phản xạ. Chuỗi và công thức phản xạ xác định cùng nghiệm liên tục từng miền.
%ID VD_3_3_2_03
\textbf{Bước 1. Trị riêng Dirichlet–Neumann.} $X(0)=0$, $X'(L)=0$. Tích phân từng phần cho $\lambda\int X^2=\int(X')^2>0$ vì nghiệm hằng thỏa $X(0)=0$ chỉ là $0$. Với $\lambda=k^2$, $X=C\sin kx$, $\cos kL=0$, nên $k_n=(n+1/2)\pi/L$, $n\ge0$.

\textbf{Bước 2. Trực giao và hệ số.} Với hai trị riêng, trừ $X_m''X_n-X_n''X_m$ và tích phân. Hạng biên $[X_m'X_n-X_n'X_m]_0^L=0$, nên tích vô hướng bằng $0$ nếu khác trị riêng. Chuẩn là $L/2$, suy $A_n=(2/L)\int f\sin k_nx$, $B_n=(2/L)\int g\sin k_nx$.

\textbf{Bước 3. Ghép nghiệm.} $u=\sum_{n\ge0}(A_n\cos ak_nt+B_n\sin ak_nt/(ak_n))\sin k_nx$. Mỗi mode đổi dấu sau $2L/a$ và có chu kỳ $4L/a$. Tính duy nhất theo năng lượng: thông lượng ở $0$ hết vì $u_t=0$, ở $L$ hết vì $u_x=0$.

\textbf{Bước 4. Dữ liệu cụ thể.} $a=1$, $L=\pi$, $f=x$, $g=\sin^3(x/2)$. Với $k_n=n+1/2$, tích phân từng phần cho $\int_0^\pi x\sin k_nx=(-1)^n/k_n^2$, nên $A_n=2(-1)^n/(\pi k_n^2)$. Từ công thức sin lập phương, $B_0=3/4$, $B_1=-1/4$:
\[u=\sum_{n\ge0}\frac{2(-1)^n}{\pi k_n^2}\cos k_nt\sin k_nx
+\frac32\sin(t/2)\sin(x/2)-\frac16\sin(3t/2)\sin(3x/2).\]
Chu kỳ cơ bản $4\pi$. $f'(\pi)=1$ trái biên Neumann, nên nghiệm không $C^1$ qua mặt sóng từ góc phải; đây là nghiệm năng lượng/từng miền, không nghiệm cổ điển toàn bao đóng.
%ID BT_3_3_3_01
\textbf{Bước 1. Công thức chung.} Hai biên cố định cho $X_n=\sin(n\pi x/L)$, $\omega_n=an\pi/L$ và hệ số $A_n=(2/L)\int fX_n$, $B_n=(2/L)\int gX_n$. Dùng $u=\sum(A_n\cos\omega_nt+B_n\sin\omega_nt/\omega_n)X_n$.

\textbf{Bước 2. Tính hệ số từng dữ liệu.}

(i) $a=L=1$. $A_1=1$, còn lại $0$. Hai lần tích phân từng phần cho
\[\int_0^1x(1-x)\sin(n\pi x)dx=\frac{2(1-(-1)^n)}{n^3\pi^3}.\]
Do đó $B_n=4(1-(-1)^n)/(n^3\pi^3)$, và
\[u=\sin\pi x\cos\pi t+\sum_{n\text{ lẻ}}\frac8{n^4\pi^4}\sin(n\pi t)\sin(n\pi x).\]
Chu kỳ $2$; dữ liệu tương thích cho $C^2$.

(ii) $a=2$, $L=\pi$, $f=\sin x$, $g=1$. $A_1=1$; $B_n=(2/\pi)\int_0^\pi\sin nx=2(1-(-1)^n)/(\pi n)$. Suy $u=\sin x\cos2t+\sum_{n\text{ lẻ}}2\sin2nt\sin nx/(\pi n^2)$, chu kỳ $\pi$. Vận tốc không tương thích ở hai đầu, nên có gãy cấp một.

(iii) $a=1$, $L=\pi$, $f=1$, $g=0$. Cùng tích phân trên cho $A_n=2(1-(-1)^n)/(\pi n)$, nên $u=\sum_{n\text{ lẻ}}4\cos nt\sin nx/(\pi n)$. Chu kỳ $2\pi$; li độ không tương thích biên tạo bước nhảy.

(iv) $a=L=1$, $f=x(1-x)/4$, $g=0$. Chia hệ số của đa thức ở (i) cho $4$: $A_n=(1-(-1)^n)/(n^3\pi^3)$. Vậy $u=\sum_{n\text{ lẻ}}2\cos(n\pi t)\sin(n\pi x)/(n^3\pi^3)$, chu kỳ $2$. $f''(0)=f''(1)=-1/2$ không khớp biên cấp hai, nên không $C^2$ qua mặt phản xạ, dù nghiệm $C^1$.

(v) Cùng $f$, $g=-1$: giữ $A_n$ trên, thêm $B_n=-2(1-(-1)^n)/(n\pi)$. Nghiệm theo công thức chung, chu kỳ $2$, không $C^1$ qua mặt sóng từ góc vì vận tốc biên khác không.

(vi) $f=0$, $g=-1/2$ trên $(0,1/2)$, $1/2$ trên $(1/2,1)$. Tách tích phân tại $1/2$:
\[B_n=-\int_0^{1/2}\sin(n\pi x)dx+\int_{1/2}^1\sin(n\pi x)dx
=\frac{2\cos(n\pi/2)-1-(-1)^n}{n\pi}.\]
Chỉ $n\equiv2\pmod4$ khác không, với $B_n=-4/(n\pi)$. Nghiệm $u=-\sum_{n\equiv2\ (4)}4\sin(n\pi t)\sin(n\pi x)/(n^2\pi^2)$ có chu kỳ cơ bản $1$. Vận tốc bước nhảy cho nghiệm liên tục nhưng gãy cấp một.

\textbf{Bước 3. Kiểm tra.} Mỗi mode thỏa PDE và biên, hai chuỗi tại $t=0$ đúng các hệ số đã tính. Chỉ đạo hàm chuỗi ở nơi hội tụ đủ; với dữ liệu không tương thích dùng phản xạ để xác định các mặt gãy, không khẳng định cổ điển toàn miền.
%ID BT_3_3_3_02
\textbf{Bước 1. Chọn hệ sine.} Trên $(0,\pi)$, $f=\min(x,\pi-x)$, $g=0$, $a=0.0002$. Tách tích phân tại $\pi/2$ và đổi biến $y=\pi-x$ ở nửa sau:
\[A_n=\frac2\pi(1+(-1)^{n+1})\int_0^{\pi/2}x\sin nx\,dx.\]
Với $n$ chẵn bằng $0$. Với $n$ lẻ, $\cos(n\pi/2)=0$, tích phân bằng $\sin(n\pi/2)/n^2$, nên chung $A_n=4\sin(n\pi/2)/(\pi n^2)$.

\textbf{Bước 2. Giải ODE.} $T_n''+a^2n^2T_n=0$, $T_n(0)=A_n$, $T_n'(0)=0$, cho
\[u=\sum_{n\ge1}\frac{4\sin(n\pi/2)}{\pi n^2}\cos(0.0002nt)\sin nx.\]
\textbf{Bước 3. Chu kỳ và tính trơn.} Mode $1$ khác không, mọi tần số là bội nguyên của $a$, nên chu kỳ cơ bản $2\pi/a=10000\pi$. Li độ có góc tại $\pi/2$; hai góc chạy theo $x=\pi/2\pm at$ rồi phản xạ ở hai đầu. Nghiệm liên tục/ năng lượng, không $C^1$ trên các nhánh truyền góc. Thế $t=0$ thu lại khai triển tam giác; biên đúng theo từng mode.
%ID BT_3_3_3_03
\textbf{Bước 1. Biên Dirichlet–Neumann.} Tách biến cho $k_n=(n+1/2)\pi/L$. Khi dịch $2L/a$, pha tăng $(2n+1)\pi$; cos và sin đều đổi dấu, nên $u(t+2L/a)=-u(t)$. Dịch thêm lần nữa cho chu kỳ $4L/a$.

\textbf{Bước 2. Kiểm tra khẳng định chu kỳ ngắn.} Nếu còn $u(t+2L/a)=u(t)$ thì $u=-u$, suy $u=0$. Lấy $f=\sin(\pi x/(2L))$, $g=0$ cho $u=f\cos(a\pi t/(2L))$ là phản ví dụ rõ ràng với chu kỳ cơ bản $4L/a$. Vì vậy chu kỳ $2L/a$ trong câu (a) không đúng cho nghiệm khác không.

\textbf{Bước 3. Biên Neumann hai đầu.} Tích phân PDE theo $x$ cho trung bình $\bar u''=(a^2/L)[u_x]_0^L=0$, nên $\bar u=\bar f+t\bar g$. Phần còn lại gồm các tần số $an\pi/L$, tất cả có chu kỳ $2L/a$. Nếu $\bar g=0$ thì nghiệm tuần hoàn; nếu $\bar g\ne0$, trung bình tăng tuyến tính và không thể tuần hoàn. Điều kiện cần và đủ là $\int_0^Lg=0$. Nghiệm hằng theo thời gian có mọi chu kỳ dương.
%ID BT_3_3_3_04
\textbf{Bước 1. Hệ cos và mode hằng.} $A_0,B_0$ là trung bình dữ liệu; mode khác không có chuẩn $L/2$. Không được dùng công thức sine của biên cố định.

\textbf{Bước 2. Câu (i).} $a=L=1$, $f=\cos2\pi x$, $g=x-1/2$. $A_2=1$, $B_0=0$. Tích phân từng phần cho
\[B_n=2\int_0^1(x-1/2)\cos(n\pi x)dx=\frac{2((-1)^n-1)}{n^2\pi^2}.\]
Do đó $u=\cos2\pi x\cos2\pi t-\sum_{n\text{ lẻ}}4\sin(n\pi t)\cos(n\pi x)/(n^3\pi^3)$. Mode $1$ khác không cho chu kỳ cơ bản $2$. $g'=1$ ở hai đầu trái đạo hàm thời gian của biên Neumann, nên không $C^2$ qua các mặt sóng góc.

\textbf{Bước 3. Câu (ii).} $a=2$, $L=\pi$, $f=\sin^2x=(1-\cos2x)/2$, $g=\sin x$. $A_0=1/2$, $A_2=-1/2$, $B_0=2/\pi$. Tích–tổng cho $B_n=2(1+(-1)^n)/(\pi(1-n^2))$ nếu $n\ne1$, và $B_1=0$ tính riêng. Suy
\[u=\tfrac12-\tfrac12\cos2x\cos4t+2t/\pi+\sum_{n\ge1}\frac{B_n}{2n}\sin2nt\cos nx.\]
\textbf{Bước 4. Kiểm tra.} Mode hằng đảm bảo vận tốc trung bình đúng; bỏ nó sẽ làm sai $g$. Nó cũng khiến câu (ii) không tuần hoàn. Biên đúng theo mỗi cos; tại $t=0$ hai chuỗi khôi phục dữ liệu. $g'(0),g'(\pi)$ khác không nên vẫn có vấn đề tương thích cấp hai.
%ID BT_3_3_3_05
\textbf{Bước 1. Trị riêng và trực giao.} $X'(0)=0$, $X(L)=0$. Tích phân từng phần cho $\lambda\int X^2=\int X'^2$. Nếu $\lambda=0$, $X$ hằng rồi $X(L)=0$ nên chỉ có nghiệm không; vì vậy $\lambda>0$. Với $\lambda=k^2$, điều kiện đầu cho $X=C\cos kx$, điều kiện cuối cho $\cos kL=0$. Suy $k_n=(n+1/2)\pi/L$. Trừ hai phương trình của $X_m,X_n$ và tích phân; các hạng biên bằng không, chứng minh trực giao.

\textbf{Bước 2. Chiếu dữ liệu và ghép.} $A_n=(2/L)\int f\cos k_nx$, $B_n=(2/L)\int g\cos k_nx$; $u=\sum(A_n\cos ak_nt+B_n\sin ak_nt/(ak_n))\cos k_nx$. Mọi mode có chu kỳ $4L/a$ và đổi dấu sau $2L/a$.

\textbf{Bước 3. Trường hợp cụ thể.} $a=1$, $L=\pi$, $f=\pi-x$, $g=\cos^3(x/2)$. Tích phân từng phần cho $\int_0^\pi(\pi-x)\cos k_nx=1/k_n^2$, nên $A_n=2/(\pi k_n^2)$. Công thức cos lập phương cho $B_0=3/4$, $B_1=1/4$. Vì vậy
\[u=\sum_{n\ge0}\frac2{\pi k_n^2}\cos k_nt\cos k_nx
+\tfrac32\sin(t/2)\cos(x/2)+\tfrac16\sin(3t/2)\cos(3x/2).\]
\textbf{Bước 4. Kết luận.} Mode đầu khác không nên chu kỳ cơ bản $4\pi$. $f'(0)=-1$ không tương thích Neumann, nên nghiệm không $C^1$ qua mặt phản xạ từ góc trái. Hai chuỗi vẫn đúng dữ liệu trong miền và từng mode thỏa PDE/biên; tính duy nhất trong lớp năng lượng theo năng lượng của hiệu.
%ID VD_3_4_2_01
\textbf{Bước 1. Hệ sine và đáp ứng nguồn cos.} $a=L=1$, $X_n=\sin(n\pi x)$, $\omega_n=n\pi$. Khai triển $x$ có hệ số $c_n=2\int_0^1x\sin(n\pi x)dx=2(-1)^{n+1}/(n\pi)$. Đặt
\[C_\omega(\alpha,t)=\frac{\cos\alpha t-\cos\omega t}{\omega^2-\alpha^2},\qquad
C_\omega(\omega,t)=\frac{t\sin\omega t}{2\omega}.\]
Đây là nghiệm $C''+\omega^2C=\cos\alpha t$, $C(0)=C'(0)=0$; nhánh cộng hưởng lấy giới hạn hoặc tích phân Duhamel.

\textbf{Bước 2. Câu (i), biên không và nguồn $x\cos\alpha t$.} Dữ liệu đầu không, nên $y_n=c_nC_{\omega_n}$ và $u=\sum_{n\ge1}c_nC_{n\pi}\sin(n\pi x)$. Nếu $\alpha=m\pi$, chỉ mode $m$ dùng nhánh cộng hưởng, không chia $0$.

\textbf{Bước 3. Câu (ii), nguồn không và biên phải $\cos\alpha t$.} Lấy $b=x\cos\alpha t$, $w=u-b$. Khi đó $Q=\alpha^2x\cos\alpha t$, $w(x,0)=-x$, $w_t(x,0)=0$. Suy
\[u=x\cos\alpha t+\sum_{n\ge1}c_n[-\cos(n\pi t)+\alpha^2C_{n\pi}(\alpha,t)]\sin(n\pi x).\]
\textbf{Bước 4. Kiểm tra và góc.} Hàm nâng cho đúng biên, ODE cho nguồn của $w$ rồi triệt tiêu với $b_{tt}$. Tại $t=0$ khai triển $-x$ bù li độ của $b$ trong miền. Nhưng ở $(1,0)$, biên phải bằng $1$ trái li độ đầu $0$, nên câu (ii) không có nghiệm liên tục trên toàn bao đóng. Câu (i) có nguồn không khớp gia tốc biên tại góc phải nếu đòi $C^2$ đến góc; chuỗi vẫn giải trong miền với cách hiểu thích hợp.
%ID VD_3_4_2_02
\textbf{Bước 1. Tính duy nhất.} Với hiệu $v$ của hai nghiệm cùng dữ liệu và biên Neumann, $v_{tt}=a^2v_{xx}$, $v_x(0,t)=v_x(L,t)=0$. Đạo hàm $E=\frac12\int(v_t^2+a^2v_x^2)$ cho $E'=a^2[v_tv_x]_0^L=0$. $E(0)=0$ suy $v_t=v_x=0$; li độ đầu bằng $0$ xác định hằng số là $0$.

\textbf{Bước 2. Nâng biên tổng quát.} Với $u_x(0,t)=\mu(t)$, $u_x(L,t)=\nu(t)$, chọn $b=\mu x+(\nu-\mu)x^2/(2L)$, vì $b_x(0)=\mu$, $b_x(L)=\nu$. Đặt $w=u-b$; nguồn mới là $Q=F+a^2(\nu-\mu)/L-b_{tt}$, hai dữ liệu mới là $f-b_0$, $g-b_{t,0}$.

\textbf{Bước 3. Tách mode.} Với $X_0=1$, $X_n=\cos(n\pi x/L)$, lấy hệ số $A_0,B_0,Q_0$ là trung bình và các hệ số $n\ge1$ là $2/L$ lần tích phân với $X_n$. Mode hằng giải $y_0''=Q_0$, nên $y_0=A_0+B_0t+\int_0^t(t-s)Q_0(s)ds$. Mode còn lại giải
\[y_n=A_n\cos\omega_nt+\frac{B_n}{\omega_n}\sin\omega_nt+\frac1{\omega_n}\int_0^t\sin(\omega_n(t-s))Q_n(s)ds.\]
\textbf{Bước 4. Trở lại nghiệm.} $u=b+y_0+\sum_{n\ge1}y_nX_n$. Kiểm tra $t=0$ sau khi cộng hàm nâng; kiểm tra từng đạo hàm biên. Công thức mode hằng khác các mode dao động, nên không áp dụng mẫu $1/\omega_0$. Nếu dữ liệu không tương thích $f'(0)=\mu(0)$, $g'(0)=\mu'(0)$ (và tương tự đầu phải), nêu rõ mức trơn và xét nghiệm yếu thay vì khẳng định cổ điển đến góc.
%ID BT_3_4_3_01
\textbf{Bước 1. Tính hệ số và nhân đáp ứng.} Trên $(0,\pi)$, khai triển $1$ theo sine có $c_n=2(1-(-1)^n)/(\pi n)$. Đặt
\[R_\omega(\alpha,t)=\frac{\sin\alpha t-(\alpha/\omega)\sin\omega t}{\omega^2-\alpha^2}.\]
Nó thỏa $R''+\omega^2R=\sin\alpha t$, $R(0)=R'(0)=0$. Khi $\alpha=\omega$, dùng $R=\sin\omega t/(2\omega^2)-t\cos\omega t/(2\omega)$.

\textbf{Bước 2. Câu (i).} $a=1$, biên và dữ liệu đầu không, nguồn $\sin\alpha t$. Chiếu nguồn cho $y_n''+n^2y_n=c_n\sin\alpha t$, nên $u=\sum c_nR_n\sin nx$. Tại cộng hưởng $\alpha=m$, thay đúng mode $m$ bằng nhánh giới hạn; nếu $c_m=0$, mode ấy không được kích thích.

\textbf{Bước 3. Câu (ii).} $a=2$, nguồn không, cả hai biên bằng $\sin\alpha t$. Lấy $b=\sin\alpha t$, $w=u-b$, suy $Q=\alpha^2\sin\alpha t$, $w_0=0$, $w_{t,0}=-\alpha$. Hệ số vận tốc là $-\alpha c_n$, vì vậy
\[u=\sin\alpha t+\sum_{n\ge1}c_n\left[-\frac\alpha{2n}\sin2nt+\alpha^2R_{2n}(\alpha,t)\right]\sin nx.\]
Đạo hàm ban đầu của hàm nâng và chuỗi triệt tiêu trong miền. Tại góc, vận tốc biên $\alpha$ trái $g=0$, nên không có nghiệm $C^1$ trên toàn bao đóng khi $\alpha\ne0$.

\textbf{Bước 4. Câu (iii).} $a=1$, nguồn $1$, biên trái $0$, biên phải $-\pi^2/2$, dữ liệu đầu không. Chọn $P=-x^2/2$: $P_{tt}-P_{xx}=1$ và đúng biên. $w=u-P$ thuần nhất, li độ $x^2/2$, vận tốc $0$. Tích phân từng phần hai lần cho
\[A_n=\frac1\pi\int_0^\pi x^2\sin nx\,dx
=-\frac{\pi(-1)^n}{n}+\frac{2((-1)^n-1)}{\pi n^3}.\]
Suy $u=-x^2/2+\sum A_n\cos nt\sin nx$. Li độ đầu và biên phải mâu thuẫn ở góc; công thức là nghiệm từng miền/ yếu. Mỗi câu đã kiểm tra nguồn qua ODE và dữ liệu qua hệ số, không chỉ ghi chuỗi hình thức.
%ID BT_3_4_3_02
\textbf{Bước 1. Giữ mode hằng Neumann.} Hệ hàm riêng trên $(0,\pi)$ là $1,\cos nx$. Mode hằng dùng $y_0''=Q_0$; các mode $n\ge1$ dùng ODE dao động và $R_\omega$ như Bài 1.

\textbf{Bước 2. Câu (i).} Biên Neumann không, $a=1$, nguồn $\sin\alpha t$ không phụ thuộc $x$. Chỉ mode $0$ có nguồn khác không. Với dữ liệu không,
\[u=\int_0^t(t-s)\sin\alpha s\,ds=t/\alpha-\sin\alpha t/\alpha^2.\]
Thế lại cho $u_{tt}=\sin\alpha t$, $u_{xx}=0$, dữ liệu không. Nếu $\alpha=0$, lấy giới hạn $u=0$.

\textbf{Bước 3. Câu (ii).} $a=2$, $\mu=\nu=\sin\alpha t$, nguồn không. Lấy $b=x\sin\alpha t$, $w=u-b$, suy $Q=\alpha^2x\sin\alpha t$, $w_0=0$, $w_{t,0}=-\alpha x$. Khai triển $x$ theo cos có $d_0=\pi/2$, $d_n=2((-1)^n-1)/(\pi n^2)$, do tích phân từng phần.

\textbf{Bước 4. Giải từng ODE và ghép.} Mode hằng là $y_0=-d_0\sin\alpha t$ (đạo hàm đầu $-\alpha d_0$ và $y_0''=\alpha^2d_0\sin\alpha t$). Các mode khác là $y_n=d_n[-\alpha\sin2nt/(2n)+\alpha^2R_{2n}]$. Do đó
\[u=(x-\pi/2)\sin\alpha t+\sum_{n\ge1}d_n[-\alpha\sin2nt/(2n)+\alpha^2R_{2n}(\alpha,t)]\cos nx.\]
Hàm nâng cho đúng đạo hàm biên; các ODE cho đúng PDE. Tại góc, $g'=0$ trái $\mu'(0)=\nu'(0)=\alpha$, nên không $C^2$ đến góc nếu $\alpha\ne0$. Tại $\alpha=2n$ phải dùng nhánh cộng hưởng, không chia mẫu bằng không.
%ID BT_3_4_3_03
\textbf{Bước 1. Chứng minh duy nhất.} Hiệu $v$ có nguồn và dữ liệu không, $v(0,t)=v_x(L,t)=0$. Đạo hàm năng lượng $E'=a^2[v_tv_x]_0^L$: ở đầu trái $v_t=0$ vì vết $v=0$, đầu phải $v_x=0$. Suy $E=E(0)=0$, $v$ hằng cả không gian và thời gian, rồi dữ liệu đầu buộc $v=0$.

\textbf{Bước 2. Nâng biên.} Lấy $b=\mu(t)+x\nu(t)$ để $b(0,t)=\mu$, $b_x(L,t)=\nu$. $w=u-b$ có biên hỗn hợp không, nguồn $Q=F-\mu''-x\nu''$, li độ $f-\mu(0)-x\nu(0)$, vận tốc $g-\mu'(0)-x\nu'(0)$.

\textbf{Bước 3. Chiếu nguồn cùng dữ liệu.} $k_n=(n+1/2)\pi/L$, $X_n=\sin k_nx$, chuẩn $L/2$. Tính $A_n=(2/L)\int w_0X_n$, $B_n=(2/L)\int w_{t,0}X_n$, $Q_n(s)=(2/L)\int Q(x,s)X_n(x)dx$.

\textbf{Bước 4. Giải và trả về biến gốc.}
\[u=b+\sum_{n\ge0}\left[A_n\cos\omega_nt+\frac{B_n}{\omega_n}\sin\omega_nt+\int_0^t\frac{\sin(\omega_n(t-s))}{\omega_n}Q_n(s)ds\right]X_n,\quad\omega_n=ak_n.\]
Với $\mu=\nu=0$, $b=0$ và dữ liệu/nguồn chiếu chính là dữ liệu gốc. Đạo hàm hai lần của tích phân mode cho $Q_n$, còn hai dữ liệu ban đầu của nó bằng không. Cộng hàm nâng xác nhận mọi biên. Cần điều kiện tương thích ở góc nếu muốn nghiệm cổ điển, chẳng hạn $f(0)=\mu(0)$, $g(0)=\mu'(0)$, $f'(L)=\nu(0)$, $g'(L)=\nu'(0)$.
%ID BT_3_4_3_04
\textbf{Bước 1. Nhận dạng một mode.} Biên cố định trên $(0,1)$ và nguồn $\sin\pi x\sin t$ gợi $u=\sin\pi x\,T(t)$. Đạo hàm $u_{xx}=-\pi^2\sin\pi x\,T$; vì tốc độ bằng $1$, PDE trở thành $T''+\pi^2T=\sin t$, dữ liệu $T(0)=T'(0)=0$.

\textbf{Bước 2. Nghiệm riêng và thuần nhất.} Nguồn tần số $1$ khác $\pi$, nên thử $T_p=C\sin t$, được $C=1/(\pi^2-1)$. Nghiệm tổng $T=C\sin t+A\cos\pi t+B\sin\pi t$. Dữ liệu đầu cho $A=0$, $C+\pi B=0$, nên $B=-1/(\pi(\pi^2-1))$.

\textbf{Bước 3. Nghiệm.}
\[u(x,t)=\frac{\sin t-\sin\pi t/\pi}{\pi^2-1}\sin\pi x.\]
\textbf{Bước 4. Kiểm tra.} Hai biên bằng không vì $\sin\pi x=0$; $T(0)=0$ và $T'(0)=(1-1)/(\pi^2-1)=0$. Thế $T''+\pi^2T$ cho $\sin t$, nên đúng nguồn. Tính duy nhất theo năng lượng của hiệu hai nghiệm biên cố định.
